\section{Introduction}\label{section: introduction}


\subsection{Aims of this paper}

The aim of this paper is to revisit fixed-point localization techniques in equivariant topological $K$-theory and employ them for an effective description of the equivariant $K$-theory ring of the affine Grassmannian. The fibers of this $K$-theory ring are given by the cohomology of varying fixed-points. Using this viewpoint, we describe $K^{\top, 0}_{T \times \Gm^{\rot}}(\Gr_G; \C)$ as a suitable deformation to normal cone (alternatively an affine blowup algebra).


The cohomology of the same varying fixed points is known to match centers of various categories $\O$ in pure representation theory. Remarkably, all this information is packaged in a single $K$-theory ring, which we completely describe in terms of topological generators. In the joint companion work \cite{HL25}, we relate the present results to classical and quantum representation theory.



Let us now explain our main results in more detail.

\subsubsection{Localization techniques in equivariant $K$-theory via derived loop stacks.}

To set up the stage, we revisit the relationship between equivariant topological $K$-theory, equivariant periodic cyclic homology, functions on derived fixed-point schemes and singular cohomology of the varying fixed points. This stems from classical fixed-point localization techniques in equivariant topology and their reinterpretation in terms of equivariant Hochschild homology along the lines of \cite{Chen20}.


\begin{motto}
The fibers of torus-equivariant $K$-theory are given by singular cohomology of the fixed-points. In other words, equivariant $K$-theory puts the cohomologies of varying fixed points into an algebraic family. Globally, equivariant $K$-theory is given by equivariant periodic cyclic homology, which has a geometric model via functions on the derived fixed-point scheme.  
\end{motto}

Combining different techniques from the literature, we arrive at the following comprehensive statement.

\begin{theorem}\label{introtheorem: localization in equivariant K-theory}
Let $T \acts X$ be a qcqs scheme over $\C$ with an action of a torus $T$. Then we have natural multiplicative $\O(T)\llp u \rrp$-linear identifications
\begin{equation*}
    K^{\top}_{T}(X; \C) \xrightarrow{\simeq} HP_{T}(X / \C) \xrightarrow{\simeq} \RG(\Fix^{\lL}_{\frac{T}{T}}(X), \O)^{tS^1}.
\end{equation*}
Taking derived fibers over a closed point $t \in T$, this equivalence specializes to
\begin{equation*}
    K^{\top}_{T}(X; \C)_{\lL t} \xrightarrow{\simeq} H_{\sing}(\Fix_{t}(X); \C)\llp u \rrp.
\end{equation*}
If $X$ is equivariantly formal, the final identification is underived, and in particular gives an isomorphism $K^{\top, 0}_{T}(X; \C)_{t} \cong H^{2\bullet}_{\sing}(\Fix_{t}(X); \C)$ of commutative rings.
\end{theorem}
\begin{proof}
\S \ref{section: K-theory and HP and cohomology of fixed points}; Theorem \ref{theorem: fibers of topological K-theory, HP and fixed-points} and Corollary \ref{corollary: K-theory in the flat situation}. 
\end{proof}

\subsubsection{Fixed-points on affine Grassmannian and  affine Schubert varieties.}
In particular, we will be interested in the action of the extended torus $T \times \Gm^{\rot}$ on the affine Grassmannian $\Gr_G$ associated to a reductive group $G$ over $k=\C$.

We can fully compute the varying fixed-points of $(x, \zeta) \in T \times \Gm^{\rot}$ in $\Gr_G$, and in fact in any affine Schubert variety $\Gr_{\leq \mu}$ inside it. These examples are equivariantly formal -- via Theorem \ref{introtheorem: localization in equivariant K-theory} above, we can effectively compute the corresponding fibers of equivariant $K$-theory.

\begin{motto}
    The fixed points of an arbitrary element $(x, \zeta)$ in the extended torus on the affine Grassmannian $\Gr_G$ are given as an infinite union of partial flag varieties (if $\zeta$ generic) resp. finite union of affine flag varieties (if $\zeta$ root of unity) for a resonant subgroup $L \leq G$. The fixed points on an affine Schubert variety $\Gr_{\leq \mu}$ inside are given by a closed union of Borel orbits resp. closed union of Iwahori orbits.
\end{motto}

Generalizing well-known computations, we first describe the reduced fixed-points $\Fix^{\red}_{(x, \zeta)}(\Gr_G)$ on the whole affine Grassmannian. We construct an associated resonant subgroup $L = L_{[x, \zeta]} \leq G$ and show that these fixed-points are given as follows.
\begin{theorem}\label{introthm: fixed points on affine grassmannians}
The reduced fixed points $\Fix^{\red}_{(x, \zeta)}(\Gr_G)$ of $(x, \zeta)\in T \times \Gm^{\rot}$ on $\Gr_G$ are given as
\begin{enumerate}
    \item[(i)] an infinite disjoint union $\coprod_{\vartheta \in X_{\bullet}/(W_L, \bullet_{\omega})} L/P_{\vartheta}$ of finite flag varieties for $L$, if $\zeta$ is generic,
    \item[(ii)] a finite disjoint union of dilated affine flag varieties for $L^{\sc}$, namely $\coprod_{\vartheta \in X_{\bullet}/(W_L, \bullet^{\ell}_{\omega})} \lo_{\ell} L^{\sc}/\eP^{\vartheta}_{\ell, L^{\sc}}$, if $\zeta$ is a root of unity of primitive order $\ell$.
\end{enumerate}
\end{theorem}
\begin{proof}
\S \ref{section: zeta generic} and \S \ref{section: zeta root of unity}; Theorems \ref{lemma: fixed points of (c, zeta) on Gr_G -- zeta generic} and \ref{lemma: fixed points of (c, zeta) on Gr_G -- zeta root of unity}.    
\end{proof}

\begin{remark}
When $x=1$ and $\zeta$ is generic, this statement is classical \cite[\S 2]{Zhu15}. When $x=1$ and $\zeta$ is a root of unity of primitive order $\ell$, this was observed in \cite[\S 4]{RW22}, \cite[\S 2]{BBSV22}. The general description of Theorem \ref{introthm: fixed points on affine grassmannians} was missing; we expect it to be useful beyond the applications in this paper. 
\end{remark}




Carefully intersecting the above with given $\Gr_{\leq \mu}$, we arrive at the following result.
\begin{theorem}
Take any affine Schubert variety $\Gr_{\leq \mu}$.
Let $(x, \zeta) \in T \times \Gm^{\rot}$ and $L=L_{[x, \zeta]} \leq G$ the corresponding resonant subgroup. Then the reduced fixed point scheme $\Fix^{\red}_{(x, \zeta)}(\Gr_{\leq \mu})$ is given by
\begin{enumerate}
    \item[(i)] a closed union of of finite Schubert varieties for $L$, if $\zeta$ is generic,
    \item[(ii)] a closed union of Iwahori orbits on $\ell$-dilated affine flag varieties $\Fl^{\vartheta}_{\ell, L^{\sc}}$ of $L^{\sc}$, if $\zeta$ is a root of unity of primitive order $\ell$.
\end{enumerate}
In both cases, we can concretely describe which finite Schubert varieties resp. closure of Iwahori orbits actually appear.
\end{theorem}
\begin{proof}
\S \ref{section: zeta generic} and \S \ref{section: zeta root of unity}; Theorems \ref{theorem: generic zeta -- fiber of fixed points and classical schubert varieties} and \ref{theorem: fiber of fixed points and classical schubert varieties general case at root of unity}.    
\end{proof}


As we explain in the companion work \cite{HL25}, this geometric statement perfectly matches known representation-theoretical phenomena in classical and quantum categories $\O$. This was, in fact, our original motivation in \cite{HL25} for conjecturing such a description in the first place. It is remarkable how accurately the resulting comparison holds.


\subsubsection{Description of equivariant $K$-theory ring of the affine Grassmannian.}
Using fixed-point localization together with the description of fixed-points on the whole $\Gr_G$ from Theorem \ref{introthm: fixed points on affine grassmannians}, we compute the $T \times \Gm^{\rot}$-equivariant complexified $K$-theory ring of the affine Grassmannian in terms of the completed ring of functions on a certain scheme ${\eN}_{\mu_{\infty}}(\eZ, T \times (T \sslash W))$ given by a gluing of deformations to normal cones of quasi-diagonals $Z_{\ell}$ inside $T \times T \sslash W$ over each root of unity $\zeta$.

\begin{motto}
The complexified equivariant $K$-theory of the affine Grassmannian is given by the completed ring of functions on infinitely many affine blowups of $T \times \Gm^{\rot} \times T\sslash W$ over each root of unity $\zeta$, centered at the $\ell$-torsion subscheme $Z_{\ell}$ of $T \sslash W$.    
\end{motto}



Such a description was originally conjectured by Roman Bezrukavnikov; it is motivated by the analogous description of equivariant cohomology from \cite[Theorem 1]{BF07}. After properly accounting for the completion, we arrive at the following. 
\begin{theorem}\label{introtheorem: bezrukavnikov description of equivariant K-theory: sc}
Let $G$ be a simply connected reductive group over $k=\C$. Then there is a functorial ring isomorphism
\begin{equation*}
   \widehat{\O}({\eN}_{\mu_{\infty}}(\eZ, T \times (T \sslash W))) \xrightarrow{\cong} K^{\top, 0}_{T \times \Gm^{\rot}}(\Gr_G; \C).
\end{equation*}
\end{theorem}
\begin{proof}
\S \ref{section: bezrukavnikov--finkelberg description of equivariant K-theory}; Theorem \ref{theorem: equivariant K-theory as deformation to normal cone -- G sssc}.    
\end{proof}

In particular, the ring $K^{\top, 0}_{T \times \Gm^{\rot}}(\Gr; \C)$ is not finitely generated, even before completing. Nevertheless, it is quite explicit, as the following remark shows.


\begin{remark}
Unraveling Theorem \ref{introtheorem: bezrukavnikov description of equivariant K-theory: sc} in terms of commutative algebra, the $K$-theory ring of the affine Grassmannian $K^{\top, 0}_{T \times \Gm^{\rot}}(\Gr_G; \C)$ is given by the completion of
\begin{equation}\label{equation: explicit K-theory introduction remark}
    \C[t_1^{\pm 1}, \dots, t_r^{\pm 1}][s^{\pm 1}_1, \dots, s^{\pm 1}_r]^W \big[\tfrac{e_j(t_1^{\ell}, \dots, t_r^{\ell}) - e_j(s_1^{\ell}, \dots, s_r^{\ell})}{1-q^{\ell}} \mid \ell \in \N, j = 1, \dots, r\big],
\end{equation}
with respect to an explicit Schubert filtration.
%
More precisely, the uncompleted ring \eqref{equation: explicit K-theory introduction remark} naturally surjects onto the equivariant $K$-theory of each affine Schubert variety $\Gr_{\leq \mu}$; the equivariant $K$-theory of the whole $\Gr_G$ is then given by completing along the induced filtration by kernels. 

In particular, we get a concrete countable infinite set of topological generators, 
\begin{equation*}
    b_{j, \ell} := \tfrac{e_j(t_1^{\ell}, \dots, t_r^{\ell}) - e_j(s_1^{\ell}, \dots, s_r^{\ell})}{1-q^{\ell}}, \qquad j =1, \dots, r, \ \ell \in \N
\end{equation*}
which we call \textit{fundamental Bezrukavnikov classes}. 
These generators have an interpretation in terms of the tautological vector bundle and are related by Adams operations. We deduce that $K^{\top, 0}_{T \times \Gm^{\rot}}(\Gr_G; \C)$ is finitely generated as a complete topological $\lambda$-ring, see Corollary \ref{corollary: topological Adams generation}. 
\end{remark}

\begin{remark}
The above description goes against some naive analogies with cohomology: the equivariant cohomology of the affine Grassmannian with complex coefficients is a polynomial (or power-series) ring on finitely many generators.

Let us explain this apparent discrepancy. In their computation of cohomology, \cite{BF07} pass to one affine blowup over the origin. (In our terms, this is due to the fact that equivariant cohomology sees only the formal completion of equivariant $K$-theory at the origin $(1,1) \in T \times \Gm^{\rot}$.) In any case, the result of this single blowup is still coincidentally isomorphic to a polynomial ring (on shifted generators). This final coincidence breaks down once we are forced to blow up infinitely many times.
\end{remark}

\begin{remark}
We want to emphasize that many interesting classes -- such as the class of the \textit{determinant line bundle} $\eL_{\det}$ -- are only present in the completion; see Lemma \ref{lemma: necessity of Schubert completion}. In other words, the Schubert completion is important and nontrivial.
\end{remark}

\begin{remark}\label{remark: variant of the K-theory of affine Grassmannian}
Our method is quite robust. In particular, it can be used to describe
    \begin{enumerate}
        \item[(1)] the corresponding result for any reductive $G$ with respect to $T^{\sc} \times \Gm^{\rot}$, 
        \item[(2)] the result with respect to an isogenous torus $T$ to $T^{\sc}$,
        \item[(3)] the $\lop G \rtimes \Gm^{\rot}$-equivariant variant,
        \item[(4)] equivariant $K$-theory of arbitrary affine flag varieties.
    \end{enumerate}    
We address these generalizations in \S \ref{section: variants and generalizations}, see in particular Theorem \ref{theorem: equivariant K-theory as deformation to normal cone -- general case}.
\end{remark}

\begin{remark}
We believe that our method may be further adapted to describe
    \begin{enumerate}
        \item[(5)] the equivariant $K$-theory with integral coefficients.
    \end{enumerate}
We will discuss this elsewhere; one description will be given in \cite{HL25}.
\end{remark}


\begin{remark}
Even for reductive $G$, it is sometimes preferable to consider the action of $G^{\sc}$ on $\Gr_G$. This has several reasons.
\begin{itemize}
    \item The equivariant $K$-theory with respect to this action is what usually appears as centers in representation theory of quantum groups; also see \cite{CK18} for such appearance.
    \item When $G$ is not simply connected, the quotient $T \sslash W$ is often singular \cite[\S 7]{BZ00}. On the other hand, $T^{\sc} \sslash W$ is always isomorphic to the affine space $\A^d = \t \sslash W$. 
    \item The simply connected variant is compatible with the theory of Kac--Moody groups. In particular, the ample generator of the Picard group of $\Gr_G$ may not be $G$-equivariant otherwise \cite[\S 3.2]{CK18}; see \cite[Remarks 4.3 and 5.8]{YZ09} for more details.
\end{itemize}
\end{remark}

\begin{remark}
The results of this paper concern topological $K$-theory of $\Gr_G$. However, we expect them to hold algebraically on the level of homotopy-invariant algebraic $K$-theory $KH$.

In \cite{Low25}, we already answered this affirmatively in type A. Namely, we showed that for affine Schubert varieties $\Gr_{\leq \mu}$ in the $GL_n$ affine Grassmannian, $KH$ is equivariantly formal and its degree zero part agrees with topological $K$-theory:
$KH^0_{T \times \Gm^{\rot}}(\Gr_{\leq \mu}) = K^{\top, 0}_{T \times \Gm^{\rot}}(\Gr_{\leq \mu})$.
Hence the results of the present paper work algebraically in this case.
\end{remark}

\begin{remark}\label{question: decompletion of K-theory}
We view the ring of functions $\O({\eN}_{\mu_{\infty}}(\eZ, T \times (T \sslash W)))$ as a useful \emph{decompletion} of $K^{\top, 0}_{T \times \Gm^{\rot}}(\Gr_G; \C)$.
However, it is not clear to us what is its $K$-theoretical meaning -- can this decompletion be realized from algebraic $K$-theory?
\end{remark}

\subsection{Companion work}
In a companion joint work in progress \cite{HL25}, we describe the tight relationship between both decompleted resp. actual equivariant $K$-theory ring $K^{\top, 0}_{T \times \Gm^{\rot}}(\Gr_G)$, and the centers of a suitable integral quantum group resp. its quantum category $\O$. We precisely pin down this quantum group. The results of \cite{HL25} work with $\Z$-coefficients, and give a new genuine application of $K$-theoretical techniques in representation theory.



\subsection{Plan of the paper}
After the current \S \ref{section: introduction}, this paper is structured as follows.

In \S \ref{section: groups and schemes} we recall some preliminary material on algebraic geometry, reductive groups, and affine flag varieties.

In \S \ref{section: equivariant K-theory and fixed points} we set up notations regarding equivariant $K$-theory, related invariants and their algebraic counterparts. The nonstandard content is our interpretation of equivariant localization in \S \ref{section: K-theory and HP and cohomology of fixed points}, in particular its formulation in Theorem \ref{theorem: fibers of topological K-theory, HP and fixed-points}.


In \S \ref{section: cocharacter graphs and resonance} we introduce the notion of resonance for $(x, \zeta) \in T \times \Gm^{\rot}$ and consequently the resonant subgroup $L_{[x, \zeta]} \leq G$. We then discuss centralizers and stabilizers attached to the point $(x, \zeta)$. We further define the {cocharacter graphs} $\Gamma$ inside $T \times \Gm^{\rot} \times T$ and explain how they put the resonant combinatorics into a family over $T \times \Gm^{\rot}$. There are no deep results, but the discussion in this section serves as a useful tool in the rest of the paper.


In \S \ref{section: fixed points on affine grassmannians}, we fully describe the varying fixed-points $\Fix^{\red}_{(x, \zeta)}(-)$ for the whole affine Grassmannian $\Gr_G$ in Theorems \ref{lemma: fixed points of (c, zeta) on Gr_G -- zeta generic} and \ref{lemma: fixed points of (c, zeta) on Gr_G -- zeta root of unity}, as well as their restriction to affine Schubert varieties $\Gr_{\leq \mu}$ inside it in Theorems \ref{theorem: generic zeta -- fiber of fixed points and classical schubert varieties} and \ref{theorem: fiber of fixed points and classical schubert varieties general case at root of unity}. These results are of independent interest.



In \S \ref{section: equivariant K-theory ring of the affine grassmannian} we move on to compute the equivariant $K$-theory ring $K^{\top, 0}_{T \times \Gm^{\rot}}(\Gr_G; \C)$ in terms of infinitely many affine blowups at roots of unity $\zeta \in \Gm^{\rot}$. The main computation for simply connected $G$ is done in Theorem \ref{theorem: equivariant K-theory as deformation to normal cone -- G sssc}. Generalizations of this result are then covered in Theorem \ref{theorem: equivariant K-theory as deformation to normal cone -- general case}.

Appendix \ref{appendix: some explicit formulas in type A} provides some concrete localization formulas in $GL_n$ and discussion of the determinant line bundle.


\subsection{Acknowledgements}

\subsubsection{Acknowledgements.}
I would like to thank to Tamás Hausel for discussions regarding both this project and our joint ongoing work relating equivariant $K$-theory to centers of quantum groups.
At the same time, I would like to thank to Roman Bezrukavnikov for sharing and discussing his conjectural expectations on equivariant topological $K$-theory of the affine Grassmannian.

Moreover, I would like to thank to Vasily Krylov for pointing out a gap in an earlier version related to appearance of pseudo-Levi subgroups in computations of fixed points, Anton Mellit for suggesting a description of the relevant $K$-theory classes in terms of the tautological vector bundle and Adams operations, Catharina Stroppel for explaining to us at an early stage of this project the appearance of cohomology rings of classical flag varieties and finite Schubert varieties as endomorphism algebras in the classical category $\O$.
%
Finally, in connection to the companion work on quantum groups, I would like to thank Quan Situ for discussing extensions of his work on quantum groups and centers of deformed categories $\O$, and Josef Svoboda for helping me notice the original independent appearance of hybrid quantum groups in quantum topology.

\subsubsection{Funding.}
This work is a revised version of the results covered in the third part of author's PhD thesis \cite[Chapter III]{LowPhD} at the Institute of Science and Technology Austria (ISTA). The results were obtained between February 2025 -- February 2026; further clarifications were implemented until September 2026. The author was supported by the DOC Fellowship of the Austrian Academy of Sciences and by the ERC grant ViaFiPoS, No 101199663.

\subsubsection{Declaration of the use of generative AI.}
The author has not used any AI tools in the writing of this article. All mathematical content was obtained independently, based on author's expertise, discussions with his colleagues and available literature.